% PCC – Justificativa
\section{Justificativa}
A justificativa para este trabalho reside na escassez de confiabilidade dos contratos em arquiteturas distribuídas. Embora a especificação OpenAPI seja um padrão altamente utilizado para documentação, podem existir divergências entre o que é declarado no contrato e o que a implementação de fato suporta, fenômeno conhecido como Contract Drift. Testes manuais e unitários tendem a validar o comportamento esperado, mas raramente exploram a negação do contrato, como entradas que deveriam ser rejeitadas, mas que, por falhas de sanitização ou lógica, causam exceções não tratadas (5xx).

A relevância técnica deste projeto concentra-se na automação do Teste Negativo Sistemático. Ao contrário de ferramentas de fuzzing genéricas que operam por força-bruta, este trabalho propõe que a inteligência do teste resida na semântica do schema. Ao transformar restrições estáticas (como minLength, pattern ou enum) em estratégias de mutação dinâmica, é possível identificar falhas de robustez que passariam despercebidas em fluxos de QA convencionais.

Portanto, o apelo do presente trabalho é transformar o contrato OpenAPI de um documento passivo em um oráculo ativo de testes, permitindo que equipes de desenvolvimento garantam a conformidade rigorosa de suas APIs sem o custo proibitivo de escrever milhares de casos de teste negativos manualmente.